\section{Conclusion}
\label{sec:conclusion}

Context compaction is the invisible infrastructure of long-running LLM
conversations --- and it is fundamentally lossy. This paper quantifies that
loss across three complementary experiments.

\textbf{Recall calibration} (\S\ref{sec:calibration}) established that recall measurement itself is
non-trivial. A questions-per-prompt effect (up to 20pp, reversing under compaction), a category-dependent recall hierarchy,
and a systematic grep-LLM gap mean that naive benchmarks conflate measurement
artifacts with real strategy differences. Any compaction evaluation that does
not control for these confounds risks measuring noise.

\textbf{Controlled single-pass compaction} (\S\ref{sec:controlled}--\S\ref{sec:singlepass}) isolated the cost of one
compaction pass. Even 5\% compaction destroys 6--7pp of recall. At 50\%, the
compacted zone is near-dead (0--7\% recall) despite 82--93\% keyword survival ---
the starkest illustration that information \textit{presence} does not equal
information \textit{retrievability}. The attention dilution effect (remaining-zone
degradation of 20pp) shows that compaction damages even untouched context.
Cross-model testing with Sonnet 4.6 --- which has no Lost-in-the-Middle effect
and 92.5\% baseline recall --- confirms that severe compaction destroys
information irrespective of model capability, while stronger models resist
light compaction significantly better ($-2.5$pp vs $-12.5$pp at C1).

\textbf{Multi-pass strategy comparison} (\S\ref{sec:strategy}) ran four strategies on a single 5M
conversation at constant fact density ($\delta=0.04$), evaluated at 5 mid-feed
checkpoints with 4--6 replicates per cell to estimate variance. The hierarchy
is \textbf{S4 FrozenRanked $>$ S3 Frozen $>$ S2 Incremental $>$ S1 Brutal} at every
checkpoint in the means. With replicates we also measured a \textbf{substantial
run-to-run variance} that is not visible in a single-shot benchmark --- the
compaction phase itself is non-deterministic at temperature 0, producing
different summaries across sessions and recall ranges spanning factor 14$\times$
on identical conversations. A \textit{single} measurement of ``strategy A vs
strategy B'' is thus untrustworthy; replicates are mandatory.

The means tell a coherent story: S4 reaches $\sim$28\% recall at 500K and decays
to $\sim$5\% at 5M; S1 stays under 10\% at every scale. The gap between S4 and S3
narrows as conversations grow, because at 5M the frozen chain is long enough
that hierarchical merging no longer differentiates the two strategies in
recall --- they converge to similar attention-dilution-limited regimes.

A complementary regression analysis on the \S\ref{sec:calibration}+\S\ref{sec:singlepass} dataset (4{,}812
observations) confirms the underlying drivers: batch size $Q$, compaction
percent, position in context, and category each act independently in
log-odds space (logistic AUC=0.84; gradient-boosted AUC=0.64, indicating
no hidden non-linearities). Compaction is monotonically destructive
($-0.054$/percent) and the spatial \textit{Lost-in-the-Middle} signature is robust.

The overarching finding is that \textbf{the bottleneck is attention, not
compression}. Keywords survive summarization; the LLM simply cannot find
them. Frozen summaries prevent the JPEG cascade but accumulate blocks that
dilute attention. This points toward a paradigm shift: from optimizing \textit{how}
we compress conversations to optimizing \textit{how} we structure preserved
information for retrieval --- indexed fact stores, RAG-augmented compaction,
and structured memory systems that bypass the attention bottleneck entirely.
